\documentclass[journal]{IEEEtran}

% Course template packages, plus the packages used by the existing report.
\usepackage{graphicx}
\usepackage[margin=1in]{geometry}
\usepackage{amsmath,amsthm,amssymb}
\usepackage{amsfonts}
\usepackage{float}
\usepackage{booktabs}
\usepackage[T1]{fontenc}
\usepackage[utf8]{inputenc}
\usepackage{enumitem}
\usepackage[hidelinks]{hyperref}
\usepackage{xcolor}

\newcommand{\system}{\textsc{ProvGate}}
\newcommand{\attack}{\textsc{Attack}}
\newcommand{\placebo}{\textsc{Placebo}}
\newcommand{\clean}{\textsc{Clean}}
\newcommand{\fullarm}{\textsc{Full}}
\newcommand{\caparm}{\textsc{Capability-Only}}
\newcommand{\promptcaparm}{\textsc{Prompt+Capability}}
\newcommand{\allowarm}{\textsc{Allow-All}}

\hyphenation{op-tical net-works semi-conduc-tor Prove-nance}

\begin{document}
\pagenumbering{arabic}
\def\arraystretch{1.2}

\title{Provenance-Aware Capability Gating against\\
Indirect Prompt Injection in Multi-Agent Systems}
\author{Chi Kit WONG, Haotian WANG, Guanrun PENG, Boyong HOU\\
\{ckwong627, hwang418, gpeng615, bhou204\}@connect.hkust-gz.edu.cn}
\date{Summer 2026}

\maketitle

\begin{abstract}
Tool-using agents can confuse data supplied by an attacker with authority
granted by the user.  We study the smallest auditable version of this problem:
ordinary capability gating checks whether a value is on an allow-list, whereas
PACT (Provenance-Aware Capability Tracking/Control) also checks whether that
value is trusted for its parameter role.  The strengthened evaluation records
six paired calls through the local \texttt{PACTToolGateway -> ToolExecutor ->
MockWorld.send\_email} boundary and eight additional policy-matrix rows.  An
externally supplied but allow-listed Bob is sent by capability-only
(\texttt{outbox\_count=1}) but is denied by PACT before \texttt{ToolExecutor}
(no tool event and \texttt{outbox\_count=0}); user-selected Bob and external
email content bound only to \texttt{content} are sent by both policies.  Exact
registered \texttt{NormalizeEmailAddress} transformations are accepted only
when source/output hashes, transform name, and trusted source authority match.
The public AgentDojo benchmark is retained only as an external attack-realism
baseline; it does not directly evaluate PACT.  The result is a small,
reproducible local mechanism demonstration rather than a claim of production or
semantic information-flow security.
\end{abstract}

\IEEEpeerreviewmaketitle

\section{Introduction}

Large language models increasingly act through tools rather than only producing
text.  Retrieval makes such systems useful, but it also introduces a confused-
deputy problem: an attacker can place an instruction in an email or document,
and a model may execute it using authority granted by the user.  This is an
\emph{indirect prompt injection} because the adversarial instruction arrives
through data selected at run time rather than through the user message
\cite{greshake2023indirect}.  Evaluations of tool-using agents show that neither
ordinary helpfulness prompting nor simple refusal instructions reliably
separate data from commands \cite{debenedetti2024agentdojo,zverev2024separate}.

This project asks:

\begin{quote}
Can a provenance-aware role check prevent an externally supplied allow-listed
value from laundering authority into a high-impact parameter, while still
allowing benign user values, low-risk external content, and explicitly
registered transformations?
\end{quote}

We deliberately make a narrow claim.  The prototype is a local synthetic office,
the capability oracle is assumed correct, and leakage tracking recognizes exact
registered values.  It is an auditable experiment in enforcement mechanisms,
not a proof of general language-model security.

\section{Background \& Motivation}

Prompt-level instruction hierarchy attempts to teach a model that trusted system
and user instructions outrank untrusted content \cite{wallace2024hierarchy}.
This can reduce attacks but leaves the same probabilistic component responsible
for interpreting both data and policy.  Capability systems instead provide an
unforgeable, least-authority description of the operations a component may
perform \cite{dennis1966capabilities,saltzer1975protection}.  Information-flow
control tracks where information came from and constrains where it may go
\cite{denning1976lattice}.

\system{} combines these ideas at the tool boundary.  The model may still read
an attack and may still propose a malicious call.  Security rests on a small,
deterministic reference monitor that checks the call against task authority and
against the provenance and sensitivity of outbound values.  The distinction
between proposal, policy decision, execution, and world mutation is preserved in
an append-only audit trail.

\section{Target System Description}

\subsection{Agents and mock tools}

The Reader Agent may call \texttt{search\_emails} and
\texttt{read\_email}.  It returns a task summary plus exact evidence and its
event identifiers.  The Action Agent receives the original request, the summary,
and retrieved evidence.  Depending on the scenario, it may propose
\texttt{read\_file}, \texttt{search\_calendar},
\texttt{send\_email}, or \texttt{create\_calendar\_event}.  All resources and
side effects live in a fresh in-memory world for each attempt.

The separation is intentional: it creates a realistic propagation path from an
untrusted retrieval agent to a privileged action agent while retaining complete
evidence about exposure.  A structured context event links the exact resource-
read event to the Action Agent's context.

\subsection{Capabilities}

Each task supplies a finite capability set.  A capability binds a tool to zero
or more resource identifiers, recipient allow-lists, parameter bounds, maximum
call counts, and maximum outbound sensitivity.  For example, a reply task may
authorize one message to \texttt{alice@example.com} while excluding added carbon
copy recipients.  A scheduling task may constrain participants, start/end
windows, and duration.  Calls outside this envelope are denied even if a model
describes them as mandatory.

\subsection{Provenance and sensitivity}

Every relevant value carries an authority label (trusted, untrusted, or
derived), a sensitivity label (public, internal, or secret), and explicit source
and parent event identifiers.  Derived values conservatively retain their
sources.  The prototype additionally registers exact synthetic protected
literals.  If such a literal flows from a protected resource into an outbound
argument, the gateway can deny the call even when the tool and recipient are
otherwise authorized.  The strengthened PACT experiment below isolates this
source-and-role check from ordinary allow-list checking.  Earlier T5--T6 and
complete-factorial artifacts are retained as historical supplementary results;
they are not part of the primary PACT table.

\section{Threat Model}

\paragraph{Adversary.}
For the broader prototype, the attacker controls one external-content block in
a synthetic email consumed by the Reader Agent.  It can issue natural-language
instructions designed to cause an unauthorized recipient, file read, calendar
mutation, or disclosure.  The primary PACT artifact does not run a live model: it uses six deterministic
policy-paired calls with fixed provenance labels through the local tool boundary,
plus an eight-row decision-only strategy matrix.  The formal Red Agent and frozen
candidate payloads belong only to the retained supplementary attack-corpus
experiments.

\paragraph{Trusted computing base.}
The scenario manifest, capability gateway, provenance tracker, \texttt{ToolExecutor},
\texttt{MockWorld}, append-only artifact writer, and independent evaluator are
trusted. In particular, the tracker and gateway enforce source-authority and
value-hash consistency before role checks; model-proposed provenance labels and
tool arguments are untrusted.  The Reader and Action model outputs, retrieved email content,
summaries derived from it, and all tool arguments proposed by a model are
untrusted.  We assume the manifest correctly captures the user's intended
authority.

\paragraph{Attacker goals.}
Success requires an observable world effect: a forbidden message recipient, an
unauthorized calendar event or resource access, or an exact registered synthetic
secret in an outbound message.  Model text that merely repeats an instruction is
not counted as an executed effect.  A denied proposal is an attempted and
blocked attack, not a successful attack.

\paragraph{Out of scope.}
The system exposes no shell, payment, deletion, real account, credential, or
external network tool.  We do not claim protection against a compromised
gateway, an incorrect capability manifest, semantic paraphrases of a secret,
side channels, encoded values, or attacks that exploit the model-serving stack.

\section{Attack \& Defense Design}

\subsection{Attacks}

The strengthened experiment has three deterministic call scenarios, each run
once under \texttt{capability\_only} and once under PACT with a fresh world.  All
six records invoke the same local \texttt{send\_email} tool when the gateway
allows the call.  The attacker can replace a high-impact recipient with a value
from an external email, or can provide ordinary text that should remain usable
as message content.  A benign user-selected Bob case checks that PACT does not
block an explicitly authorized value.

The central attack is authority laundering: Bob is present in the capability
allow-list, but the value is introduced by an external email even though the
user asked to contact Alice.  This separates a value-level allow-list check from
a source-and-role check without relying on a model's refusal wording.

\begin{table}[t]
\centering
\scriptsize
\caption{Six paired calls through the real local tool boundary.}
\label{tab:pact-e2e}
\resizebox{\columnwidth}{!}{%
\begin{tabular}{llcccc}
\toprule
Case & Policy & $C$ & $P$ & Tool event & Outbox \\
\midrule
External Bob, allow-listed recipient & Capability-only & Allow & Deny & Yes & 1 \\
External Bob, allow-listed recipient & PACT & Allow & Deny & No & 0 \\
User-selected Bob & Capability-only & Allow & Allow & Yes & 1 \\
User-selected Bob & PACT & Allow & Allow & Yes & 1 \\
External text, user-selected Alice & Capability-only & Allow & Allow & Yes & 1 \\
External text, user-selected Alice & PACT & Allow & Allow & Yes & 1 \\
\bottomrule
\end{tabular}}
\end{table}

The first pair is the authority-laundering case: capability-only reaches
\texttt{MockWorld.send\_email}, whereas PACT denies before \texttt{ToolExecutor}.
The third pair binds external data only to the low-risk \texttt{content} role;
the user still supplies the Alice recipient.

\subsection{Defenses}

Ordinary capability-only checks compare each proposed value against the
capability allow-list but ignore its provenance.  PACT performs that same value
check and then requires trusted provenance for the high-impact roles
\texttt{recipient}, \texttt{target}, and \texttt{control}.  External values
remain permitted for the low-risk \texttt{content} role.

The transformation registry is an exact-match registry: the source-value hash,
output-value hash, transform name, and trusted source authority must all match one
registered entry.  The strengthened matrix uses the semantically valid
\texttt{NormalizeEmailAddress} transform.  A user-sourced normalized address is
accepted; an unregistered transform is denied.  For the external-source row, the
same source/output hashes are registered under trusted USER authority, but the
EXTERNAL source metadata does not match that exact key and is denied.  This is
provenance verification, not fuzzy string matching.

\section{Evaluation}

\subsection{Experimental Setup}

The primary artifact is \texttt{artifacts/pact-strengthened-v3/}.  It contains
six end-to-end records (three cases paired across \texttt{capability\_only} and
PACT), an eight-row strategy matrix, a canonical JSONL decision log, an
architecture file, a strategy plot, and a SHA-256 manifest.  Every end-to-end
row starts from a fresh \texttt{MockWorld}; the gateway calls the existing
\texttt{PACTToolGateway -> ToolExecutor -> MockWorld.send\_email} path.  A PACT
deny is therefore observed as no tool event and no outbox mutation, rather than
as a decision-only counter.

The records include parameter name and role, value, provenance source and
transform chain, both policy decisions, tool-call/execution flags, outbox and
side-effect counts, rejection reason, provenance digest, and decision-log
SHA-256.  The JSON contains six case-policy records, while the CSV expands their
arguments into 18 parameter rows.  The eight strategy rows are deliberately
decision-only and exercise recipient, low-risk content, and control (the
representative high-trust role) bindings plus registered, unregistered, and
externally sourced \texttt{NormalizeEmailAddress} metadata.  Although
\texttt{target} is also declared high-trust in the policy, it is not separately
instantiated in this small matrix.

The reproducible entry point is
\texttt{scripts/}\texttt{run\_pact\_strengthened.py} with output directory
\texttt{artifacts/}\texttt{pact-strengthened-v3}.  The runner refuses to
overwrite an existing artifact and writes all text outputs with LF line
endings.

\begin{figure}[t]
\centering
\fbox{%
\begin{minipage}{0.86\columnwidth}
\centering
\texttt{user / external email}\par
$\downarrow$ attach immutable provenance\par
\texttt{recipient | target | content | control}\par
$\downarrow$ exact capability + PACT checks\par
\texttt{ALLOW: ToolExecutor -> MockWorld.send\_email}\par
\texttt{DENY: no tool event + no outbox side effect}
\end{minipage}}
\caption{PACT enforcement at the local tool boundary.}
\label{fig:pact-flow}
\end{figure}

\paragraph{Decision-log excerpt.}
A representative blocked row records
\texttt{case\_id=external-bob-allowlisted}, \texttt{role=recipient}, and
\texttt{source=external:email-evil}.  It records
\texttt{capability=Allow}, \texttt{pact=Deny}, \texttt{executed=0},
\texttt{tool\_events=[]}, and \texttt{outbox\_count=0}; the complete row also
stores the provenance digest and a deterministic decision-record SHA-256 in the
committed JSONL artifact.  This digest is computed over the canonical decision
payload; it is not a literal hash of the JSONL line containing the record.

\subsection{Outcome Definitions}

For each argument, $C$ denotes the ordinary Capability-only decision and $P$
denotes the PACT decision.  $X$ denotes an executed tool call and $S$ denotes
a local MockWorld side effect.  The mechanism checks require $X=S=1$ only when
$C=P=1$.  A denied call must have $X=S=0$.  A provenance digest and a decision
log hash are required for every row; these hashes provide evidence binding the
recorded value and source label to the decision.

\subsection{Results}

The six real-tool records reproduce the central contrast shown in
Table~\ref{tab:pact-e2e}.  Under capability-only, the external allow-listed Bob
call is allowed and produces one tool event plus one outbox message.  Under
PACT, the same call has $C=\mathrm{Allow}$ but $P=\mathrm{Deny}$; it is blocked
before \texttt{ToolExecutor}, with no tool event and zero outbox entries.  User
Bob and external content sent to the user's Alice recipient are both allowed by
both policies and each produces one outbox message.

The eight-row strategy matrix is:

\begin{table}[t]
\centering
\scriptsize
\caption{Strengthened PACT policy matrix (decision-only rows).}
\label{tab:pact-strategy}
\begin{tabular}{p{0.50\columnwidth}cc}
\toprule
Strategy row & Capability-only & PACT \\
\midrule
User $\rightarrow$ recipient, no transform & Allow & Allow \\
External $\rightarrow$ recipient, no transform & Allow & Deny \\
External $\rightarrow$ content, no transform & Allow & Allow \\
User $\rightarrow$ recipient, registered NormalizeEmailAddress & Allow & Allow \\
User $\rightarrow$ recipient, unregistered transform & Allow & Deny \\
External $\rightarrow$ recipient, registered name but EXTERNAL source & Allow & Deny \\
User $\rightarrow$ control, no transform & Allow & Allow \\
External $\rightarrow$ control, no transform & Allow & Deny \\
\bottomrule
\end{tabular}
\end{table}

All eight rows pass the capability allow-list.  PACT allows the user recipient,
external content, registered user transformation, and user control rows (4/8),
while denying the externally sourced high-trust bindings and the unverified
transform claims (4/8).  Because the matrix is decision-only, its rows do not
claim a tool side effect; the side-effect claim comes only from the six paired
records.  The machine-readable files are \texttt{e2e\_results.csv/json},
\texttt{strategy\_results.csv/json}, and \texttt{decision\_log.jsonl}; their
hashes are recorded in \texttt{manifest.json}.

\subsection{External AgentDojo Baseline}

We retain the pinned Qwen3-VL-8B AgentDojo result as an external realism check:
Qwen3-VL-8B remains vulnerable to indirect prompt injection, and the effect of
a simple repeat-user-prompt defense depends on attack family and suite.  This
result motivates the problem but does not evaluate PACT and is not merged with
Table~\ref{tab:pact-e2e}.  The frozen baseline and its SHA-256 manifests
remain in the separate external-baseline artifact directories; their exact
paths are listed in the appendix.

\section{Analysis and Discussion}

\subsection{Interpretation}

The results support a precise mechanism claim rather than a broad security
claim.  Ordinary capability gating answers ``is this value allowed?''  PACT adds
``is this source trusted for this role?''  The strengthened real-tool pair shows
that this distinction changes an observable side effect: an allow-listed Bob
from an external email is sent by capability-only but is stopped before
\texttt{ToolExecutor} by PACT.  PACT still allows external text at
\texttt{content}, preserves an explicitly selected Bob, and accepts only a
hash-verified registered \texttt{NormalizeEmailAddress} transformation.

\subsection{Limitations}

The capability manifest is assumed correct and the strengthened experiment uses
an in-memory local \texttt{MockWorld}, not a production messaging service.
The tool boundary is real within the repository, but all accounts, messages,
and side effects are synthetic.  Provenance is explicit and hash-based, not
semantic: unregistered decoding, paraphrases, partial values, and
transformations omitted from the exact-match registry are outside the claim.
The experiment uses one tool contract, one local implementation, one
deterministic policy, and no multi-seed or multi-model variance.  The eight-row
matrix is decision-only; only the six paired calls measure actual outbox side
effects.  The AgentDojo benchmark is retained as an external attack-realism
baseline only and is not merged with PACT metrics.  Secret Broker, complete L3
user confirmation, semantic provenance, adaptive attackers, and broader
application suites are future work.

\subsection{Ethics and Safety}

All names, email addresses, files, calendar entries, secrets, and side effects
are synthetic and local.  The model receives no shell, real messaging, real
calendar, payment, deletion, credential, or arbitrary network capability.
Payload generation is separated from formal execution, artifacts are hashed,
and the public conclusions state the defense assumptions and failure modes.  The
work aims to improve defensive measurement.  Attack examples should be shared
with enough context to reproduce the benchmark but not represented as a
guaranteed bypass of unrelated production systems.

\subsection{Reproducibility}

The current primary artifact is self-contained under
\texttt{artifacts/pact-strengthened-v3/}.  Its
\texttt{e2e\_results.json} contains six case-policy records, while
\texttt{e2e\_results.csv} expands their arguments into 18 parameter rows;
\texttt{strategy\_results.csv/json} contain eight decision-only rows;
\texttt{decision\_log.jsonl} contains canonical structured decisions; each
\texttt{decision\_log\_sha256} is a deterministic digest of its decision
record payload, not a literal JSONL line hash;
\texttt{strategy\_matrix.svg} and \texttt{architecture.mmd} provide compact
visual summaries; and \texttt{manifest.json} records the code commit plus
SHA-256 for every output.  The manifest's \texttt{outputs} and
\texttt{output\_sha256} maps are identical.  All persisted text is LF-canonical.

The replay command runs
\texttt{PYTHONPATH=src python} followed by
\texttt{scripts/}\texttt{run\_pact\_strengthened.py}; it writes to the
\texttt{pact-strengthened-v3} directory and refuses to overwrite an existing
artifact.  The artifact is bound to commit
\texttt{3b7f886b3e5b10fc26bf1e62ff7ffd24cd5ac0e0}.  The full commit is recorded in the manifest.  The earlier
\texttt{pact-minimum-v2}/v1 artifacts and pinned AgentDojo outputs are retained
as historical or external-baseline evidence and are not merged into the PACT
metrics.

\subsection{AI-Use Disclosure}

Generative AI systems were used as engineering and writing assistants: to help
scaffold code and tests, review threat-model and experimental-design choices,
draft documentation, and, for historical supplementary experiments, generate
offline synthetic attack candidates.  Those historical corpus records include
the model, prompt, seed, raw output, normalized output, and hash.  The primary strengthened PACT artifact is deterministic and does not
depend on a model response; its six paired calls execute through the local tool
boundary.  Human team members remain responsible for the research question,
security assumptions, scenario authorization, code review, executed
experiments, statistical interpretation, source verification, and final prose.
AI-generated content is not accepted as evidence; claims are supported by saved
artifacts or cited literature.

\section{Peer Evaluation}

\subsection{Individual Contribution Statements}

\begin{itemize}[nosep]
  \item Chi Kit WONG: Led project scoping and security-experiment design; developed the threat model and provenance-aware capability-gating architecture; designed and froze the original and hardened offline attack corpora; specified the $2 \times 2$ ablation and deterministic sink-pressure tests; coordinated experiment execution and artifact validation; maintained the GitHub/Overleaf reproducibility package and integrated the final report.
  \item Haotian WANG:
  \item Guanrun PENG:
  \item Boyong HOU:
\end{itemize}

\paragraph{Equal contribution.}
All four members made equal contributions.  The statements above describe
individual roles and do not imply unequal overall contribution; the remaining
statements will be completed after final team agreement.

\subsection{Teammate Evaluation Scores \& Justifications}

The team reports equal contribution for all four members.  No differentiated
scores or justifications are asserted in this report; any separate course
peer-evaluation form should use the same equal-contribution statement.

\section{Conclusion}

We presented PACT, a provenance-aware capability monitor for multi-agent tool
calls, and strengthened its evidence with a real local tool boundary.  In six
paired calls through \texttt{PACTToolGateway -> ToolExecutor -> MockWorld},
ordinary capability checks send an externally supplied allow-listed Bob,
whereas PACT rejects the same high-trust recipient before the executor and
leaves no outbox side effect.  User-selected Bob and external email text bound
only to \texttt{content} are sent by both policies.  An eight-row strategy matrix
further shows exact role and transformation behavior: a registered
\texttt{NormalizeEmailAddress} result is accepted, while unregistered or
externally claimed transforms are denied.  The artifact contains the records,
decision logs, provenance digests, LF-canonical outputs, and manifest hashes
needed to reproduce these claims.  The result is a focused course-project
mechanism demonstration on a synthetic local system, not a claim of complete
production or semantic information-flow security.

\nocite{*}
\bibliographystyle{IEEEtran}
\bibliography{references}

\newpage
\section*{Appendices}

\subsection*{Artifact and replay pointers}

The primary PACT artifact is under
\path{artifacts/pact-strengthened-v3/}.  The six real-tool records are in
\path{e2e_results.csv} and \path{e2e_results.json}; the eight strategy rows are
in \path{strategy_results.csv} and \path{strategy_results.json}; the decision
log is \path{decision_log.jsonl}; and the plot, architecture, and SHA-256
manifest are \path{strategy_matrix.svg}, \path{architecture.mmd}, and
\path{manifest.json}.  The artifact is bound to the strengthened runner commit
and uses LF-canonical output.  Historical \path{pact-minimum-v2}/v1 artifacts,
frozen attack corpora, factorial runs, and AgentDojo summaries remain under
their existing paths as supplementary or external-baseline evidence and are not
mixed with the strengthened PACT metrics.  The report can be rebuilt with
\path{bash scripts/compile_report.sh}; deterministic controller tests use the
repository's \texttt{test} Conda environment.

\end{document}
